\section{Discussion}
\label{sec:discussion}

\paragraph{Two things that are easy to conflate}
Our experiments pull apart two claims that usually travel together. The first is about a \emph{class} of topology.
Connectome-scale sparse signed recurrence, trained only to predict a video model's latent, ends up holding a state
that path-integrates: heading and position are linearly readable through 8 to 20 dark steps, and a GRU trained on the
same objective, at either parameter budget, never forms this code, even when we look for it with a non-linear probe.
We find this strong, because the GRU predicts the latent better. The code is not a by-product of prediction quality;
if anything it trades against it. The second claim is about the \emph{identity} of the real synapses, and here the
evidence is modest: about one percent of latent error, 2\degree\ to 6\degree\ of heading, reproducible across seeds,
environments, data sizes and subgraphs, but never the cause of the effect and absent from the behavioural readouts.
A reader who takes away only ``sparse signed recurrence integrates, and the fly's version is a little better'' has
understood the paper.

\paragraph{Why might sparse recurrence integrate?}
We can offer a hypothesis but not yet a proof. A sparse signed graph at this scale has a large reservoir of slow,
weakly coupled modes. The leaky update in \cref{eq:student} drives those modes with the action input, and the linear
readout is free to leave them alone because they cost almost nothing in latent error. A gated unit with a dense
$174^2$ or $762^2$ recurrence is pushed by the same objective to spend its capacity on prediction, and nothing in the
loss rewards it for keeping an integral around. The probes on random and shuffled sparse graphs fit this picture: the
code appears whenever the recurrence is sparse, signed and large. The obvious next experiments, a spectral analysis
of the trained $W$ and lesions of its slow modes, are left for future work.

\paragraph{Limitations}
We list them in the order we think reviewers should weigh them. (i) The behavioural readouts are linear and that is
where the result stops short: homing reaches 37\,\% on the NYC twin and 10\,\% on the arena against a 92 to 100\,\%
ceiling, and goal-vector navigation does not beat random actions. A readout fit on more rollouts, or a recurrent
readout, may close part of this gap, but we have not shown it. (ii) The information ceiling is trivial because turns
are fixed-angle and actions are given in the dark. The turn-noise control in \cref{tab:homing} is a first step toward
non-trivial integration; the full experiment with continuous turn angles remains to be run. (iii) The environments are
synthetic, and on real handheld video no student beats a copy-last baseline. (iv) The teacher is ViT-B at 256\,px,
the smaller \vjepa variant. (v) The NYC twin uses one location and a random, not geographic, episode split. (vi) The
51 EPG neurons do not carry the heading code; the biological ring attractor does not emerge at this scale, and we do
not claim it.

\paragraph{Broader impact}
The connectome is released by \flywire under CC-BY-NC~4.0; our derived graph file inherits that licence and is rebuilt
by one script rather than redistributed. No human data are used. If there is a lesson for the field, it is a
methodological one: biological wiring is a source of structural priors to be tested with matched controls, not a
blueprint to be copied, and the controls are where most of the information is.

\section{Conclusion}
A video foundation model can be distilled into a world model whose recurrence is the fly's wiring diagram. The wiring
makes it a worse predictor and a better integrator: the state keeps a linearly readable heading and position through
the dark that a GRU trained on the same objective never forms, with a small, consistent advantage for the real
synapses over shuffled ones. Turning that code into reliable blind behaviour is the open problem, and we release
everything needed to work on it.
